\section{Robustness}\label{sec:robust}

\subsection{Dependence on 2022}

The CPI strategy earned most of its profit in 2022, when inflation dominated markets and every CPI
print moved the expected path of the policy rate (Figure~\ref{fig:byyear}): 79\% of its net P\&L
in \ES\ and 74\% in \NQ. Excluding 2022, its Sharpe ratio falls from 1.03 to 0.58 in \ES\ and from 1.00
to 0.68 in \NQ\ for the naive rule. \YM\ repeats the pattern: 2022 contributes 4.1 of its 6.8
percentage points, and its Sharpe ratio excluding 2022 is about 0.5. The broad book, which spreads its risk over about ten releases,
is less concentrated: its Sharpe ratio excluding 2022 is 0.28 in \ES, 0.44 in \NQ\ and 0.40 in \YM, and the
surprise books are positive in both the validation and the test years of the equal-terms
comparison. We read the result as state dependence at the level of the macro regime: the drift is
largest when macro data drive policy expectations.

\begin{figure}[htbp]
  \centering
  \includegraphics[width=\linewidth]{fig18_cpi_pnl_by_year_4markets}
  \caption{Net P\&L of the CPI strategies by calendar year (\% of notional), \ES, \NQ, \YM\ and \ZN.}
  \label{fig:byyear}
\end{figure}

\subsection{Entry delay and holding period}

The \ES\ CPI Sharpe ratio remains between 0.55 and 0.71 when the entry is delayed by one, two or five
minutes with a 30-minute hold, and between 0.51 and 0.62 in \NQ. \YM\ is more sensitive to the
delay: from 0.84 with no delay to 0.54, 0.33 and 0.17 after one, two and five minutes. Short holds are fragile: a
five-minute hold loses most of its Sharpe ratio when the entry is delayed by one minute. The edge is a
slow continuation over about half an hour, not a fast reaction.

\subsection{More windows per release}

A natural way to trade more often is to trade several consecutive windows after each release.
Figure~\ref{fig:ladder} shows that only the first window is clearly profitable, and mainly for the
surprise: the surprise's rank correlation with the return falls from 0.139 in the first 30 minutes
in \ES\ (0.088 in \NQ\ and 0.091 in \YM) to between $-0.06$ and $+0.05$ in the next four windows.
The analogues have no information in any window in \ES, \NQ\ and \ZN; in \YM\ their first window
earns 0.47 and the later ones lose. Every additional window adds turnover and cost but not profit; five-window
books lose between 29\% and 48\% of notional in \ES\ over 2019--2026, and on \YM\ the five-window
surprise and DTW books have Sharpe ratios of $-0.34$ and $-0.23$.

\begin{figure}[htbp]
  \centering
  \includegraphics[width=\linewidth]{fig28_ladder_sharpe_by_window_4markets}
  \caption{Net Sharpe ratio of the surprise and DTW signals traded in each of the five 30-minute
  windows after a release, \ES, \NQ, \YM\ and \ZN, 2019--2026.}
  \label{fig:ladder}
\end{figure}

\subsection{Construction choices and multiple testing}

Pattern-based results are sensitive to reasonable construction choices. For the \ES\ DTW ladder,
the Sharpe ratio on 2021--2026 at a quarter tick ranges from 0.17 to 0.64 across the number of
neighbours, the width of the trading gate and whether its thresholds are fixed or re-estimated;
across 63 variations of waiting time, path length and neighbour weights, it ranges from $-0.45$ to
0.61. The original implementation of the ladder, with its own family grouping and data, reaches
0.64 in \ES\ on 2021--2026 at a quarter tick (random-sign $p=0.04$) but 0.11 in \NQ, and its
direction forecast has a negative rank correlation with returns in the training years. The best of
many such variations overstates what an investor could expect
\citep{Harvey2016, Bailey2014}. The surprise approach has few free choices, and its main
constants were fixed before the test years were scored. The pattern-based books also pass a leak
test: extending the price path five minutes past the entry raises the \ES\ DTW Sharpe ratio from
0.64 to 2.86, so the timing of the real book is correct and the information is not coming from
the future. In the state-dependence analysis, analytic HC3 and permutation $p$-values agree closely,
and we report only effects that survive the false-discovery correction of \citet{Benjamini1995}.

\subsection{Limitations}

Our results are subject to five limitations. First, costs are modeled rather than observed; fills
in the first minute after a release may be worse than two ticks. Second, the sample contains about
120 releases per family, and CPI-only results rest on 50--90 trades. Third, the event-game setting was
chosen on the full sample, and the \YM\ results were produced after the \ES\ and \NQ\ results were
known, although with unchanged code and constants. Fourth, the test years have been examined more than once during the
project, so they are not an untouched hold-out for any approach. Fifth, for near-zero CPI surprises
the sign depends on which CPI line is used: the headline and the family average disagree on 20 of
86 releases, which changes the CPI Sharpe ratio from 1.03 to 0.51 in \ES.
